\section{Modelo de Datos y Persistencia}
\label{sec:database}

El sistema utiliza una base de datos relacional \textbf{PostgreSQL} alojada en \textbf{Supabase} para mantener el estado de las conversaciones y registrar los eventos históricos para su análisis estadístico.

\subsection{Esquema de Tablas}
La base de datos consta de dos tablas principales: \texttt{whatsapp\_sessions} y \texttt{validation\_logs}.

\subsubsection{Tabla: whatsapp\_sessions}
Esta tabla almacena el estado actual de la sesión del usuario. Se indexa por el número de teléfono del cliente para garantizar transiciones de estado rápidas.

\begin{table}[ht]
\centering
\caption{Atributos de la tabla whatsapp\_sessions}
\footnotesize
\renewcommand{\arraystretch}{1.4}
\begin{tabular}{p{1.8cm} p{2cm} p{3cm}}
\toprule
\textbf{Columna} & \textbf{Tipo} & \textbf{Descripción} \\
\midrule
\texttt{phone} & TEXT (PK) & Número de teléfono del usuario. \\
\texttt{state} & TEXT & Estado actual de la máquina de estados. \\
\texttt{front\_result} & JSONB & Metadatos técnicos de la cara frontal. \\
\texttt{back\_result} & JSONB & Metadatos técnicos de la cara trasera. \\
\texttt{updated\_at} & TIMESTAMPTZ & Fecha y hora del último mensaje. \\
\bottomrule
\end{tabular}
\end{table}

\subsubsection{Tabla: validation\_logs}
Esta tabla registra cada una de las interacciones y transacciones procesadas por el bot. Es la fuente de datos principal para el Dashboard de monitoreo.

\begin{table}[ht]
\centering
\caption{Atributos de la tabla validation\_logs}
\footnotesize
\renewcommand{\arraystretch}{1.4}
\begin{tabular}{p{2.2cm} p{1.8cm} p{2.8cm}}
\toprule
\textbf{Columna} & \textbf{Tipo} & \textbf{Descripción} \\
\midrule
\texttt{id} & SERIAL (PK) & Identificador único correlativo. \\
\texttt{phone} & TEXT & Número de teléfono del usuario. \\
\texttt{timestamp} & TIMESTAMP\-TZ & Fecha y hora del procesamiento. \\
\texttt{node} & TEXT & Nombre del nodo ejecutado. \\
\texttt{whatsapp\_\allowbreak message\_id} & TEXT & Identificador wamid de Meta. \\
\texttt{lambda\_\allowbreak score} & NUMERIC & Score de fraude obtenido. \\
\texttt{lambda\_\allowbreak confidence} & NUMERIC & Confianza de detección de cédula. \\
\texttt{lambda\_\allowbreak response} & JSONB & Respuesta JSON completa de AWS. \\
\texttt{latency\_\allowbreak aws\_ms} & INTEGER & Latencia del API de AWS (ms). \\
\texttt{latency\_\allowbreak gemini\_ms} & INTEGER & Latencia del API de Gemini (ms). \\
\texttt{status} & TEXT & Estado del resultado de la petición. \\
\texttt{is\_photocopy} & BOOLEAN & Indicador binario de fraude. \\
\bottomrule
\end{tabular}
\end{table}

\subsection{Restricciones de Pooling en Supabase}
Dado que Supabase implementa un Connection Pooler (PgBouncer) configurado en \textbf{modo transacción} para manejar altas concurrencias (puerto \texttt{6543}), se configuró explícitamente el parámetro \texttt{statement\_cache\_size=0} en la librería \texttt{asyncpg} en el backend. 

Esta configuración es obligatoria para evitar errores de ejecución en la base de datos debido al reuso indebido de prepared statements a nivel de sesión en el pooler.
